\section{Main concepts}
\begin{figure*}[ht!]
    \begin{subfigure}{.19\textwidth}
        \centering
        \includegraphics[height=.11\textheight]{decomposition-3D-no-decomp}
        \caption{}
        \label{fig:decomp-no-decomp}
    \end{subfigure}
    \hfill
    \begin{subfigure}{.19\textwidth}
        \centering
        \includegraphics[height=.11\textheight]{decomposition-3D-3-cell}
        \caption{}
        \label{fig:decomp-volume}
    \end{subfigure}
    \hfill
    \begin{subfigure}{.19\textwidth}
        \centering
        \includegraphics[height=.11\textheight]{decomposition-3D-2-cell}
        \caption{}
        \label{fig:decomp-face}
        \end{subfigure}
    \hfill
    \begin{subfigure}{.19\textwidth}
        \centering
        \includegraphics[height=.11\textheight]{decomposition-3D-1-cell}
        \caption{}
        \label{fig:decomp-edge}
    \end{subfigure}
    \hfill
    \begin{subfigure}{.19\textwidth}
        \centering
        \includegraphics[height=.11\textheight]{decomposition-3D-0-cell}
        \caption{}
        \label{fig:decomp-vertex}
    \end{subfigure}
    \caption{Cell decomposition of a geometric 3D object}
    \label{fig:cell-decompositions}
\end{figure*}

\begin{figure*}[ht!]
  \begin{subfigure}{.19\textwidth}
    \centering
    \includegraphics[height=.11\textheight]{orbit-3D-vertex-123.pdf}
    \caption{}
    \label{fig:orbit-3D-vertex-123}
  \end{subfigure}
  \hfill
  \begin{subfigure}{.19\textwidth}
    \centering
    \includegraphics[height=.11\textheight]{orbit-3D-edge-023.pdf}
    \caption{}
    \label{fig:orbit-3D-edge-023}
  \end{subfigure}
  \hfill
  \begin{subfigure}{.19\textwidth}
    \centering
    \includegraphics[height=.11\textheight]{orbit-3D-face-013}
    \caption{}
    \label{fig:orbit-3D-face-013}
  \end{subfigure}
  \hfill
  \begin{subfigure}{.19\textwidth}
    \centering
    \includegraphics[height=.11\textheight]{orbit-3D-volume-012}
    \caption{}
    \label{fig:orbit-3D-volume-012}
  \end{subfigure}
  \hfill
  \begin{subfigure}{.19\textwidth}
    \centering
    \includegraphics[height=.11\textheight]{orbit-3D-edge-02}
    \caption{}
    \label{fig:orbit-3D-edge-02}
  \end{subfigure}

  \caption{Orbit decomposition of a geometric 3D object}
  \label{fig:house-orbit-decomposition}
\end{figure*}

In this section, we present the generalized maps, graph transformation rules and their subsequent concepts which are necessary to the understanding of our contribution.

\subsection{Generalized maps}\label{ssec:g-maps}

\emph{Generalized maps} (or G-maps) \cite{lienhardt1991topological,
  Damiand-Lienhardt14} allow the representation of manifold geometric objects
(with or without boundaries), based on a cellular $n$-dimensional topological
structure. %
The representation of an object as a G-map comes intuitively from its
decomposition into topological cells (vertices, edges, faces, volu\-mes, and so
on). %
For example, the 3D topological object (Fig.~\ref{fig:decomp-no-decomp}) can be
decomposed into two volumes (Fig.~\ref{fig:decomp-volume}): a cube and a pyramid. %
These volumes are \emph{linked} along their common faces with a $3$-link, drawn
in green. %
The index "$3$" means that the link connects two $3$-dimensional (possibly a
single one) volumes. %
In the same way, volumes are split into faces connected with blue $2$-links
(Fig.~\ref{fig:decomp-face}). %
Then, faces are split into edges connected with red $1$-links
(Fig.~\ref{fig:decomp-edge}). %
Lastly, edges themselves are split into vertices with black $0$-links
(Fig.~\ref{fig:decomp-vertex}) to produce the $3$-G-map describing the objects
shown in Fig.~\ref{fig:decomp-no-decomp}. %
A G-map is therefore a graph, the nodes (named \emph{darts}) are vertices of edges 
of faces of volumes and the arcs are $i$-links. %
By convention, border darts have $3$ loops which are not represented to %preserve the figures' readability. % 
make the figures easier to read.

G-maps have conditions guaranteeing objects consistency, for example, two faces
are always linked along an edge.

Topological cells are not explicitly represented in G-maps but only implicitly
defined as subgraphs named \emph{orbits}. %
They can be computed using graph traversals defined by an originating dart and
by a given set of link labels. %
For example, the 0-cell (or the object's vertex) incident to some dart $a$
(Fig.~\ref{fig:orbit-3D-vertex-123}) is the subgraph which contains $a$
and all darts reachable from $a$, using links labelled by $1$, $2$ or $3$ and
the links themselves. %
% (\textit{i.e.} darts $a$, $b$, $c$, $d$, $e$, $f$, $g$ and $h$) and the links
% themselves. %
This subgraph is denoted by $G\orbit{1,2,3}(a)$ where $\orbit{1,2,3}$ is the
\emph{type} of the orbit and models a vertex. %
% the face $BCED$ (Fig.~\ref{fig:house-ABCDE}). %
The $1$-cell %adjacent 
(or edge) incident to $a$ 
(Fig.~\ref{fig:orbit-3D-edge-023}) is the subgraph $G\orbit{0,2,3}(a)$ which
contains $a$ and all the reachable darts using links labelled by $0$, $2$ or
$3$ % (\textit{i.e.} darts $a, b, l$ and $m$) %
and the corresponding links. %
% It represents the topological edge $BC$. %
% Thanks to $2$-loops on the object border, $G\orbit{0, 2}(c)$
% (Fig.~\ref{fig:house-1-cell-c}) also represents the edge $CE$ matching dart
% $c$. %
The $2$-cell (or face) incident to $a$ (Fig.~\ref{fig:orbit-3D-face-013})
is the orbit $G\orbit{0,1,3}(a)$. %
The $3$-cell (or volume) incident to $a$ (Fig.~\ref{fig:orbit-3D-volume-012}) is the
orbit $G\orbit{0,1,2}(a)$. %
% The $0$-cell adjacent to $a$ (or the object vertex matching dart $a$)
% (Fig.~\ref{fig:house-0-cell-a}) is the subgraph $G\orbit{1, 2}(a)$ and
% represents vertex $B$. %
Note that orbits are more general than cells. %
For example, the vo\-lu\-me edge $G\orbit{0,2}(a)$ (Fig.~\ref{fig:orbit-3D-edge-02})
is the $\orbit{0,2}$-orbit incident to $a$.


\subsection{Graph transformation rules}
\label{ssec:rules}

Jerboas's \cite{belhaouari2014jerboa} graph transformation rules allow the
formalization of operations over G-maps. %

\begin{figure*}[ht!]
  \centering
  \includegraphics[height=38pt]{rule-face-extrusion}
  \caption{Extrusion rule}
  \label{fig:rule-face-extrusion}
\end{figure*}

\begin{figure}[ht!]
  \centering
  \begin{subfigure}[t]{.45\columnwidth}
    \centering
    \includegraphics[width=.6\textwidth]{demo-fig-extrusion-left-match.pdf}
    \caption{Graph $G$}
    \label{fig:demo-fig-extrusion-left-match}
  \end{subfigure}
  \hfill
  \begin{subfigure}[t]{.45\columnwidth}
    \centering
    \includegraphics[height=.6\textwidth]{demo-fig-extrusion-right-match.pdf}
    \caption{Graph $H$}
    \label{fig:demo-fig-extrusion-right-match}
  \end{subfigure}
  \caption{Face's extrusion (rule Fig.~\ref{fig:rule-face-extrusion})}
  % \label{fig:before-after-triangulation}
  \label{fig:before-after-extrusion}
\end{figure}

In a few words, a \emph{rule} $r : L \longrightarrow R$ and a \emph{match}
$m : L \to G$ to a G-map $G$, describe the transformation
$G \longrightarrow^{r,m} H$ from $G$ to $H$. %
The match $m$ allows the replacement of a subgraph of $G$ described by the
left-hand side of the rule $L$ with another one described by the right-hand side
$R$, in order to produce $H$. %

Informally, in the extrusion rule (Fig.~\ref{fig:rule-face-extrusion}), the
left-hand side is made of only one node $n_{1}$ (orange) labelled with the $\orbit{0,1}$ face
type: this way, it can match any face. %
%
For the match $m : n1 \mapsto 6$ from $L$ to $G$
(Fig.~\ref{fig:demo-fig-extrusion-left-match}), the node $n_{1}$ matches 
the whole face $G\orbit{0, 1}(6)$. %
%
On the right side, the node $n_{1}$ label remains $\orbit{0,1}$. %
This means that, %when 
after applying the rule, the matched face $\orbit{0,1}(6)$ % is
has been 
preserved, in other words $G\orbit{0,1}(6) = H\orbit{0,1}(6)$ (Fig.~\ref{fig:before-after-extrusion}). %
%
In $R$, the new node $n_{2}$ (blue) creates, %in $H$, 
a copy of the matched face
in $H$.
However, $n_{2}$'s label is $\orbit{0, \_}$ meaning that $0$-links are preserved
and $1$-links are deleted. %
Therefore, $n_{2}$ creates face edges $\orbit{0}$ from the edges of the matched face. %
In a similar way, $n_3$ (pink) creates a\-no\-ther copy of the matched face. %
Because it is labelled $\orbit{\_, 2}$, $0$-links are deleted, $1$-links
relabelled with $2$, $n_3$ creates edge vertices $\orbit{2}$ from the matched
face's vertices. %
Finally, the nodes $n_{4}$, $n_{5}$, $n_{6}$ create the same orbits than nodes
$n_{3}$, $n_{2}$ and $n_{1}$, respectively. %
The nodes' labels, called \emph{implicit arcs}, match the highlighted links 
(Fig.~\ref{fig:demo-fig-extrusion-right-match}). %

The arc between $n_1$ and $n_2$, called \emph{explicit arc}, is $2$-labelled in the extrusion rule
(Fig.\ref{fig:rule-face-extrusion}) and $2$-links one-to-one the preserved orange
darts and the created blue darts (Fig. \ref{fig:before-after-extrusion}). %
Similarly, the explicit arc between $n_2$ and $n_3$, $1$-links one-to-one the blue and pink darts. %

The node $n_{1}$ (Fig.~\ref{fig:rule-face-extrusion}) is a \emph{preserved node}
because it belongs to both the left and right-hand sides of the rule. %
Nodes $n_{2}$ to $n_{6}$ are \emph{created nodes} because they belong only to
the right-hand side. %
\emph{Deleted nodes} belong only to the left-hand side. %
Note that the extrusion rule does not have any deleted node. %

The \emph{orbit} notion is extended to patterns of rules. %

% For example, in the right-hand side of the extrusion rule
% (Fig.~\ref{fig:rule-face-extrusion}), the $\orbit{0,1}$-orbit (a face) incident
% to node $n_{2}$ contains the nodes $n_{2}$, $n_{3}$, $n_{4}$ and $n_{5}$ and the
% explicit $0$- and $1$-links that connect them. %
% Additionally, an $\orbit{o}$-orbit is said to be \emph{complete} if each node in
% the orbit matches one link per label of $\orbit{o}$ either explicitly or
% implicitly. %
% For example, the $\orbit{0,1}$-orbit incident to node $n_{2}$ is complete since
% each of its nodes has both $0$- and $1$-links either explicitly or implicitly. %
% Note that an $\orbit{o}$-orbit in a graph is entirely matched by a rule pattern
% if and only if the corresponding $\orbit{o}$-orbit in the pattern is complete. %
% For example, the $\orbit{0,1}$-orbit incident to the dart $8$ in the graph $H$ 
% (Fig.~\ref{fig:demo-fig-extrusion-right-match}) is completely matched by the 
% extrusion rule because the $\orbit{0,1}$-orbit incident to $n_{2}$ is complete. % 
%
%For example, in the right-hand side of the triangulation rule,
%the $\orbit{0, 2}$-orbit (an edge) incident
%to node $n_{0}$ contains the single node $n_0$, and the $\orbit{0, 2}$-orbit
%incident to $n_1$ contain the two nodes $n_1$ and $n_2$ and the explicit
%$0$-link which connects them. %
%Moreover, an $\orbit{o}$-orbit is said to be \emph{complete} if each node in
%the orbit matches one link per label of $\orbit{o}$ either explicitly or
%implicitly. %
%For example, the $\orbit{0, 2}$-orbit incident to node $n_{0}$ in the right-hand
%side of the triangulation rule (Fig.~\ref{fig:rule-triangulation}) is not
%complete, because node $n_0$ has no $2$-link, either implicitly or explicitly. %
%Conversely, the $\orbit{0, 2}$-orbit incident to $n_1$ is complete, because its
%two nodes $n_0$ and $n_1$ are incident to the explicit $0$-link and both have an
%implicit $2$-link. %
%Note that an $\orbit{o}$-orbit in a graph is entirely matched by a rule pattern
%if and only if the corresponding $\orbit{o}$-orbit in the pattern is complete. %
%For example, the $\orbit{0, 2}$-orbit incident to $a$ in the graph $I$
%(Fig.~\ref{fig:demo-fig-triangulation-right-match}) is partially matched by the
%triangulation rule (only dart $a_{0}$ and its $0$-linked neighbour are matched),
%because the $\orbit{0, 2}$-orbit incident to $n_0$ in the rule is not
%complete. %
%Conversely, the $\orbit{0, 2}$-orbit incident to $a_1$ in graph $I$ is entirely
%matched by the triangulation rule (both darts $a_1$ and $h_1$ are matched by
%node $n_1$ and darts $a_2$ and $h_2$ are matched by node $n_2$), because the
%$\orbit{0, 2}$-orbit incident to $n_1$ in the rule is complete. %

Jerboa's rules provide syntactic properties %guaranteeing G-maps consistency preservation 
which guarantee the preservation of the consistency of G-maps \cite{belhaouari2014jerboa}. %

\subsection{Orbit tracking}

\begin{figure}[ht!]
    \centering
    \includegraphics[height=45pt]{rule-face-triangulation-3D-boxed.pdf}
    \caption{Face triangulation rule}
    \label{fig:rule-triangulation}
\end{figure}

\begin{figure}[ht!]
\begin{subfigure}{.45\columnwidth}
    \centering
    \includegraphics[width=.7\textwidth]{cube-article-view.pdf}
    \caption{Graph H}
    \label{fig:demo-fig-triangulation-left-match}
\end{subfigure}
\hfill
\begin{subfigure}{.45\columnwidth}
    \centering
    \includegraphics[width=.7\textwidth]{triangulated-cube-article-view.pdf}
    \caption{Graph I}
    \label{fig:demo-fig-triangulation-right-match}
\end{subfigure}
\caption{Front-left face's triangulation (rule
  Fig.~\ref{fig:rule-triangulation})}
\label{fig:concept-triangulation}
\end{figure}

Rules contain the necessary information characterizing the topological changes 
affecting an orbit throughout an application \cite{gaide2023automatic}. %
Thus, the tracking of an orbit is automatically made
without any addition other than the rules' syntactic analysis. %

Consider the example of the triangulation (Fig.~\ref{fig:rule-triangulation})
and its application on $H$'s dart $33$ (Fig.~\ref{fig:demo-fig-triangulation-left-match}). %
The left-hand side of the rule matches the whole green square face $H\orbit{0,1,3}(33)$ 
while its right-hand side splits it into four green triangles in graph $I$ (Fig. \ref{fig:demo-fig-triangulation-right-match}). %

Similarly to G-maps' orbits, in the left-hand side of the rule, the $\orbit{0,1,3}$-orbit incident to $n_{0}$
is the orbit containing the node rea\-cha\-ble through arcs labelled in $\orbit{0,1,3}$ and those same arcs. %
This left-hand side orbit is written $L\orbit{0,1,3}(n_{0})$. %
Therefore, $L\orbit{0,1,3}(n_{0})$ is the orbit matching the green face. %

In green 
(Fig.~\ref{fig:rule-triangulation}), an $1$-arc connects the nodes $n_{0}$ and $n_{1}$ and a $0$-arc
connects $n_{1}$ and $n_{2}$, thus forming the  $\orbit{0,1,3}$-orbit (face) incident to $n_{0}$. %
This orbit matches the four faces resulting from the % rule's application
application of the rule 
(Fig.~\ref{fig:demo-fig-triangulation-right-match}). %
The syntactic analysis of the rule allows us to deduce that the face orbit is
\emph{split} along its implicit $1$-arcs because the second implicit $1$-arc of $n_0$ in
the left-hand side is relabelled outside of the face $\orbit{0, 1, 3}$-orbit type
in the right-hand side. %
Consequently, the matched face of graph $H$ is split along its vertices'
$1$-links and into four faces in graph $I$. %
Similarly, a rule \emph{merges} two or more $\orbit{o}$-orbits when a $k$-th implicit arc is relabelled from $i$ to $j$,
where $i\not\in\orbit{o}$ and $j\in\orbit{o}$, while there was no such $k$-th implicit arc
in any node of the left-hand side. %

In red, 
the vertex orbit $R\orbit{1,2,3}(n_{2})$ incident to $n_{2}$, matches the vertex $I\orbit{1,2,3}(127)$.
Since $R\orbit{1,2,3}(n_{2})$ only contains $n_{2}$ which %is 
has been 
created, thus the orbit 
itself 
is \emph{created} and the % rule's application 
application of the rule 
creates the vertex $I\orbit{1,2,3}(127)$. % 
%Moreover, because the implicit arcs $1,2,3$ of $n_2$ are relabelled from the matched implicit links
%$0,1,3$ of $n_0$, this created vertex $\orbit{1,2,3}$-orbit comes from the matched face $\orbit{0,1,3}$-orbit
%incident to the node $n_0$ in the left-hand side of the rule. %
%Consequently, the corresponding orbit (red) in the graph $I$ is created with the
%matched face as its \emph{origin}. %


In blue, 
the face vertex $\orbit{1,2}$-orbit incident to the preserved node
$n_{0}$ matches the face vertices of the 
green face such as the vertex orbit $H\orbit{1,2}(33)$.
The node $n_{1}$ is added to the orbit through a $1$-arc, $R\orbit{1,2}(n_{0})$, 
thus modifying it. %
Consequently, the matched face's vertices of graph $H$ are modified in graph $I$. %

Finally, the face edge $\orbit{0}$-orbit incident to the preserved node $n_{0}$
matches the green face's edges in $H$. %
Since no node is added nor deleted from the orbit nor any arc is relabelled, thus 
the orbit is \emph{not modified}.

%With such 
Through this
analysis, topological changes can be logged within bulletin boards
which are automatically computed without re\-qui\-ring any other
intervention.

%%% Local Variables:
%%% mode: latex
%%% TeX-master: "../main"
%%% End:
